% Pista ana-24s-q1 · Anàlisi
% Procedència: PAU Catalunya 2024 · Convocatòria de setembre · Sèrie 3 · Qüestió 1
\documentclass[11pt,a4paper]{article}
\usepackage{pista}
\pistainfo{Catalunya 2024}{Convocatòria de setembre}{Sèrie 3}

\begin{document}

\section*{\textcolor{blauPista}{Qüestió 1}}

\noindent Considereu la funció polinòmica $f(x) = 3x^{13} + 5x^{3} + 2$.

\subsection*{\textit{a)} \normalfont Justifiqueu que la seva gràfica talla l'eix de les abscisses en un punt de l'interval $[-2, 0]$. Doneu un interval de longitud $0{,}5$ on es trobi aquest punt de tall. \hfill\textnormal{[1,25~punts]}}

\pista{Aquí has d'aplicar el Teorema de Bolzano: si $f$ és contínua a l'interval i canvia de signe entre els extrems, llavors hi ha alguna arrel. Calcula $f(-2)$ i $f(0)$ i comprova-ho. Per a localitzar un interval de longitud $0{,}5$, prova valors intermedis ($-1{,}5$, $-1$, $-0{,}5$) fins a trobar un canvi de signe entre dos valors consecutius.}

\subsection*{\textit{b)} \normalfont Estudieu les zones de creixement i de decreixement, i els màxims i els mínims de $y = f(x)$. Quants punts de tall té exactament la gràfica d'aquesta funció amb l'eix de les abscisses? Justifiqueu la resposta. \hfill\textnormal{[1,25~punts]}}

\pista{Calcula $f'(x)$ i analitza el seu signe. Observa que $f'(x)$ és una suma de monomis de grau parell amb coeficients positius: quin signe pot tenir? Dedueix-ne el creixement i si hi ha màxims o mínims. Finalment, combina-ho amb l'apartat a) per justificar quants punts de tall té la gràfica amb l'eix d'abscisses.}

\end{document}
